\documentclass[conference]{IEEEtran}
\IEEEoverridecommandlockouts
\usepackage{cite}
\usepackage{amsmath,amssymb,amsfonts}
\usepackage{algorithmic}
\usepackage{graphicx}
\usepackage{textcomp}
\usepackage{xcolor}
\usepackage{booktabs}
\usepackage{multirow}
\usepackage{array}

% ── Bahasa Indonesia: ganti label bawaan IEEEtran ──────────────
\renewcommand{\abstractname}{Abstrak}
\renewcommand{\IEEEkeywordsname}{Kata Kunci}
\renewcommand{\refname}{Daftar Pustaka}

\def\BibTeX{{\rm B\kern-.05em{\sc i\kern-.025em b}\kern-.08em
    T\kern-.1667em\lower.7ex\hbox{E}\kern-.125emX}}

% ───────────────────────────────────────────────────────────────
\begin{document}

\title{Pengenalan Aksara Jawa Tulisan Tangan untuk Klasifikasi\\
120 Kelas Silabik Menggunakan Augmentasi Berbasis\\
Preservasi Struktur dan Weighted Loss\\
pada Dataset Tidak Seimbang}

\author{
\IEEEauthorblockN{1\textsuperscript{st} Fathan Fardian Sanum}
\IEEEauthorblockA{\textit{School of Computing} \\
\textit{Telkom University}\\
Bandung, Indonesia \\
fathansanum@student.telkomuniversity.ac.id}
\and
\IEEEauthorblockN{2\textsuperscript{nd} Syahdan Rizqi Ruhendy}
\IEEEauthorblockA{\textit{School of Computing} \\
\textit{Telkom University}\\
Bandung, Indonesia \\
syahdanruhendy@student.telkomuniversity.ac.id}
\and
\IEEEauthorblockN{3\textsuperscript{rd} Ardia Alif Ramadhan}
\IEEEauthorblockA{\textit{School of Computing} \\
\textit{Telkom University}\\
Bandung, Indonesia \\
ardiaramadhan@student.telkomuniversity.ac.id}
}

\maketitle

% ───────────────────────────────────────────────────────────────
\begin{abstract}
Aksara Jawa (\textit{Hanacaraka} / \textit{Carakan}) merupakan sistem
penulisan \textit{abugida} yang menjadi warisan budaya penting
masyarakat Jawa, Indonesia.
Pengenalan otomatis aksara Jawa tulisan tangan pada cakupan
penuh 120 kelas silabik, yang dibentuk dari kombinasi 20 konsonan
dasar (\textit{nglegena}) dengan 6 tanda vokal, masih sangat
terbatas, terutama pada kondisi data yang tidak seimbang secara
alami.
Penelitian ini mengatasi celah tersebut menggunakan dataset
\textit{Indonesian Local Script Characters} yang memiliki rasio
ketidakseimbangan kelas sebesar 33,10$\times$ antara kelas mayoritas
(Vokal A, hingga 695 gambar per kelas) dan kelas minoritas
(Vokal E, I, O, U, dan E-taling, hanya sekitar 21 gambar per kelas).
Kami mengusulkan \textit{pipeline} persiapan data empat langkah:
(1) stratified split 70:15:15 per kelas,
(2) integrasi data variasi tulisan tangan,
(3) augmentasi berbasis preservasi struktur yang hanya menggunakan
transformasi fotometrik dan morfologis tanpa transformasi geometris
yang dapat mengubah makna visual aksara, dan
(4) fungsi loss berbobot (\textit{weighted cross-entropy}).
Empat arsitektur \textit{transfer learning} dievaluasi menggunakan
\textit{pipeline} yang sama: EfficientNet-B0, ResNet-18, VGG-16,
dan MobileNetV3-Large.
EfficientNet-B0 mencapai akurasi tertinggi 97,08\% dengan F1-score
makro 0,9110, sedangkan ResNet-18 menunjukkan keseimbangan
performa terbaik antara kelas mayoritas dan minoritas
(F1 Vokal A = 0,8972 vs.\ F1 Vokal non-A = 0,8976, selisih 0,0004).
Evaluasi dilaporkan secara terpisah untuk kelas mayoritas (Vokal A)
dan minoritas (Vokal non-A), memberikan gambaran yang lebih
transparan mengenai kemampuan generalisasi model pada kondisi
ketidakseimbangan nyata.
\end{abstract}

\begin{IEEEkeywords}
aksara Jawa, pengenalan karakter tulisan tangan, dataset tidak seimbang,
augmentasi preservasi struktur, \textit{transfer learning},
EfficientNet-B0, ResNet-18, VGG-16, MobileNetV3-Large,
\textit{abugida}, \textit{weighted loss}
\end{IEEEkeywords}

% ───────────────────────────────────────────────────────────────
\section{Pendahuluan}

Aksara Jawa (\textit{Hanacaraka} atau \textit{Carakan}) adalah sistem
penulisan tradisional yang telah digunakan selama berabad-abad dalam
naskah sastra, dokumen keagamaan, dan catatan resmi di Jawa,
Indonesia~\cite{ifriza2026}.
Meskipun memiliki nilai sejarah dan budaya yang tinggi, tingkat
penguasaan dan penggunaan aksara Jawa terus menurun akibat dominasi
aksara Latin dalam sistem pendidikan modern~\cite{faizin2025,wisesty2026}.
Kondisi ini menimbulkan ancaman nyata terhadap kelestarian warisan
budaya tersebut.

Salah satu strategi pelestarian yang efektif adalah digitalisasi
naskah melalui \textit{Optical Character Recognition} (OCR),
yang memungkinkan konversi dokumen historis ke format digital yang
dapat dicari dan diakses untuk berbagai aplikasi seperti mesin
penerjemah, platform pembelajaran digital, dan basis data warisan
budaya.
Namun, pengembangan sistem OCR yang akurat untuk aksara Jawa
menghadirkan tantangan teknis tersendiri yang berakar pada
kompleksitas struktur tulisannya.

Aksara Jawa diklasifikasikan sebagai \textit{abugida}
(alphasillabari), yaitu sistem tulisan di mana setiap simbol dasar
merepresentasikan konsonan dengan vokal inheren, sedangkan vokal
lain ditandai melalui diakritik tambahan yang disebut
\textit{sandhangan}~\cite{wicaksono2024,widiarti2026}.
Kombinasi 20 konsonan dasar dengan 6 posisi vokal (Vokal A, E, I, O,
U, dan E-taling) menghasilkan \textbf{120 kelas silabik yang berbeda}.
Widiarti dan Adji~\cite{widiarti2026} mencatat bahwa sistem aksara
Jawa secara lengkap, termasuk bentuk \textit{pasangan}, dapat
menghasilkan lebih dari 11.000 unit silabik yang berbeda,
menegaskan kompleksitas inheren sistem ini untuk keperluan OCR.

Pengenalan pada cakupan 120 kelas menghadapi dua hambatan teknis
utama.
Pertama, beberapa pasangan karakter memiliki kemiripan visual
yang sangat tinggi.
Nindya et al.~\cite{nindya2025} mengidentifikasi pasangan ``ha''/``la''
dengan \textit{Character Likeness Rate} (CLR) sebesar 25,3\%
(keduanya merupakan cerminan horizontal yang hampir identik),
serta pasangan ``sa''/``da'' (CLR 15,8\%).
Kedua, satu-satunya dataset publik yang mencakup seluruh 120 kelas
silabik, yaitu \textit{Indonesian Local Script Characters}~%
\cite{wisesty2026,sulistiyo2025}, memiliki distribusi kelas yang
sangat tidak seimbang: kelas Vokal A rata-rata memiliki $\approx$500
gambar per kelas, sementara 100 kelas lainnya (Vokal E, I, O, U,
dan E-taling) hanya memiliki $\approx$21 gambar per kelas,
menghasilkan rasio ketidakseimbangan sebesar 33,10$\times$.

Penelitian ini memberikan tiga kontribusi utama:
\begin{enumerate}
  \item \textbf{Metodologis:}
    Sebuah \textit{pipeline} data empat langkah yang menggabungkan
    stratified split, integrasi data variasi, dan strategi
    \textit{augmentasi berbasis preservasi struktur} yang secara
    eksplisit mengecualikan transformasi geometris yang terbukti
    dapat merusak struktur visual aksara~%
    \cite{wisesty2026,nindya2025}.
  \item \textbf{Empiris:}
    Evaluasi komparatif empat arsitektur \textit{transfer learning}
    (EfficientNet-B0, ResNet-18, VGG-16, MobileNetV3-Large) untuk
    pengenalan 120 kelas silabik aksara Jawa pada dataset yang tidak
    seimbang secara alami (IR\,=\,33,10$\times$), melengkapi
    studi-studi sebelumnya yang menggunakan dataset seimbang
    buatan~\cite{susantoIJAI2023,susantoIJECE2023}.
  \item \textbf{Analitis:}
    Pelaporan metrik kinerja secara terpisah per kelompok vokal
    (Vokal A vs.\ Vokal non-A), yang mengungkap kesenjangan
    generalisasi yang tersembunyi oleh metrik agregat dan membantu
    memilih arsitektur yang tepat bergantung pada prioritas
    (akurasi tertinggi vs.\ keseimbangan antar kelas).
\end{enumerate}

Sisa makalah ini disusun sebagai berikut.
Bagian~\ref{sec:related} mengulas penelitian terkait.
Bagian~\ref{sec:dataset} mendeskripsikan dataset.
Bagian~\ref{sec:method} memaparkan metode yang diusulkan.
Bagian~\ref{sec:setup} merinci pengaturan eksperimen.
Bagian~\ref{sec:results} menyajikan dan mendiskusikan hasil.
Bagian~\ref{sec:conclusion} menutup makalah.

% ───────────────────────────────────────────────────────────────
\section{Tinjauan Pustaka}
\label{sec:related}

\subsection{Pengenalan Aksara Jawa: Studi 120 Kelas}

Susanto et al.~\cite{susantoIJECE2023} mengembangkan CNN khusus
dengan empat lapisan konvolusional dan dua lapisan \textit{fully
connected} untuk klasifikasi 120 kelas, mencapai akurasi 97,29\%
pada dataset seimbang berisi 14.400 gambar (120 per kelas) dari
30 penulis.
Kelompok yang sama~\cite{susantoIJAI2023} kemudian mengusulkan
CNN 12 lapisan dengan augmentasi selektif (transformasi afin,
proyektif, dan penskalaan minimal dengan rotasi dibatasi pada
$\pm3^{\circ}$ untuk menghindari perubahan makna aksara),
mencapai akurasi 99,65\%.
Kedua studi beroperasi pada dataset yang diseimbangkan secara
artifisial, sehingga kondisi data tidak seimbang secara alami
belum pernah dieksplorasi.
Pembatasan sudut rotasi pada~\cite{susantoIJAI2023} secara langsung
menjadi motivasi strategi augmentasi yang diusulkan dalam penelitian ini.

\subsection{Pengenalan Multi-Aksara Nusantara}

Wisesty et al.~\cite{wisesty2026} mengusulkan MNetNCR berbasis
MobileNetV3~\cite{howard2019mobilenetv3} untuk lima aksara Nusantara,
mencapai F1\,=\,0,979 untuk aksara Jawa (20 kelas dasar)
\textit{tanpa augmentasi}.
Secara kritis, penerapan augmentasi justru \textit{menurunkan}
performa dari 0,9788 menjadi 0,9579, yang dikaitkan dengan
transformasi yang mengganggu struktur visual khas aksara Jawa.
Temuan ini menjadi motivasi empiris utama bagi pembatasan
augmentasi yang diterapkan dalam penelitian ini.

Nindya et al.~\cite{nindya2025} menerapkan VGG-16~\cite{simonyan2015vgg}
pada klasifikasi multi-aksara Bali, Jawa, dan Sunda, mencapai
akurasi 93,19\% untuk subset Jawa (20 kelas).
Penelitian ini memperkenalkan metrik CLR dan mengidentifikasi
pasangan ``ha''/``la'' (CLR 25,3\%) dan ``sa''/``da'' (CLR 15,8\%)
sebagai karakter yang paling sering tertukar karena kemiripan
visualnya, mengkonfirmasi bahwa \textit{horizontal flip} harus
dihindari.

Sulistiyo et al.~\cite{sulistiyo2025} membandingkan metode
\textit{shallow learning} (LightGBM, XGBoost, Random Forest,
SVM, CatBoost) pada dataset yang sama, dengan LightGBM mencapai
F1-score tertinggi sebesar 77,15\% untuk aksara Jawa.
Hasil ini menegaskan bahwa \textit{deep learning} jauh mengungguli
pendekatan \textit{shallow learning} pada tugas klasifikasi
multi-kelas yang kompleks.

Ifriza et al.~\cite{ifriza2026} menggunakan ResNet-18~\cite{he2016resnet}
yang di-\textit{fine-tune} untuk mengenali aksara Jawa 20 kelas
dari dataset Hanacaraka (1.562 gambar), mencapai akurasi 95,53\%.
Penelitian ini menjadi satu-satunya rujukan \textit{transfer learning}
ResNet untuk aksara Jawa.

\subsection{Penanganan Ketidakseimbangan pada Aksara Jawa}

Faizin et al.~\cite{faizin2025} dari ITS mengembangkan SkelBAGAN
(\textit{Skeleton-based Balancing GAN}) untuk dataset aksara Jawa
dengan rasio ketidakseimbangan hingga 81,78$\times$.
SkelBAGAN mengintegrasikan \textit{autoencoder} kerangka karakter
dan \textit{adversarial loss} berbasis kerangka untuk menghasilkan
gambar sintetis yang mempertahankan integritas struktur aksara.
Ini merupakan satu-satunya penelitian sebelumnya yang secara eksplisit
menangani ketidakseimbangan pada pengenalan aksara Jawa, sehingga
menjadi landasan komparatif bagi pendekatan tanpa GAN yang kami
usulkan.

\subsection{OCR Aksara Abugida Asia Tenggara}

Wicaksono et al.~\cite{wicaksono2024} melakukan tinjauan literatur
sistematis tentang segmentasi kata untuk aksara abugida Asia
Tenggara (Thai, Burma, Lao, Khmer, Jawa, dan Sunda), menemukan
bahwa tidak ada penelitian dalam lima tahun terakhir yang membahas
segmentasi kata untuk aksara Jawa maupun Sunda.
Widiarti dan Adji~\cite{widiarti2026} menetapkan evaluasi baseline
untuk OCR aksara Jawa \textit{cetak} (bukan tulisan tangan),
mencatat tantangan tambahan yang ditimbulkan oleh variasi tulisan
tangan.
Kedua survei ini mengkonfirmasi bahwa pengenalan karakter terisolasi,
yaitu cakupan penelitian ini, merupakan batas penelitian terkini
untuk OCR aksara Jawa.

% ───────────────────────────────────────────────────────────────
\section{Dataset}
\label{sec:dataset}

\subsection{Sumber dan Struktur}

Penelitian ini menggunakan dataset \textit{Indonesian Local Script
Characters}, dikurasi oleh Aditya Firman Ihsan et al.\ di
Universitas Telkom dan tersedia di Mendeley Data.
Dataset ini juga digunakan pada penelitian-penelitian
sebelumnya~\cite{wisesty2026,sulistiyo2025}.
Untuk aksara Jawa, dataset terdiri dari dua sub-direktori:
\begin{itemize}
  \item \texttt{all\_class/}: 6 kelompok vokal $\times$ 20
    sub-folder konsonan = 120 direktori kelas, dengan total
    12.191 gambar.
  \item \texttt{variations/}: 8 direktori konsonan (ba, ga, ha,
    ka, la, nya, ra, ta) berisi 1.066 gambar yang merepresentasikan
    variasi gaya tulisan tangan untuk kelas Vokal A yang bersesuaian.
\end{itemize}
Gambar didominasi format JPEG (91\%) dan PNG (9\%), saluran warna
RGB (98,3\%), dengan resolusi asli $\approx362\times347$ piksel.
Seluruh gambar distandarisasi ke ukuran $224\times224$ piksel
sebagai input model.

\subsection{Distribusi Kelas dan Karakterisasi Ketidakseimbangan}

Tabel~\ref{tab:dist} merangkum distribusi per kelompok vokal.
Dua puluh kelas Vokal A masing-masing memiliki 426--695 gambar
($\mu = 499{,}8$), sedangkan 100 kelas non-A hanya memiliki
21--22 gambar per kelas.
Rasio ketidakseimbangan (maks/min) adalah \textbf{33,10$\times$},
dengan median per kelas sebesar 22, jauh di bawah rata-rata
101,59, yang mengkonfirmasi distribusi sangat condong ke kanan
dan didominasi oleh sejumlah kecil kelas mayoritas.

\begin{table}[htbp]
\caption{Distribusi Kelas Dataset per Kelompok Vokal}
\label{tab:dist}
\begin{center}
\renewcommand{\arraystretch}{1.2}
\begin{tabular}{lccc}
\toprule
\textbf{Kelompok Vokal} & \textbf{Kelas} &
  \textbf{Gambar/Kelas} & \textbf{Total} \\
\midrule
Vokal A        & 20 & 426--695 & 9.995 \\
Vokal E        & 20 & 21--22   &   440 \\
Vokal I        & 20 & 21--22   &   440 \\
Vokal O        & 20 & 21--22   &   440 \\
Vokal U        & 20 & 21--22   &   440 \\
Vokal E-taling & 20 & 21--22   &   440 \\
\midrule
\textbf{Total} & \textbf{120} & 21--695 & \textbf{12.195} \\
\bottomrule
\multicolumn{4}{l}{\footnotesize IR\,=\,33,10$\times$;
  Median\,=\,22; Rerata\,=\,101,59; Std\,=\,188,94}
\end{tabular}
\end{center}
\end{table}

Sub-direktori \texttt{variations/} (IR\,=\,2,87$\times$ dalam
subset ini) menyediakan gambar tambahan untuk 8 dari 20 konsonan
Vokal A, digunakan pada Langkah 2 dari \textit{pipeline} yang
diusulkan.

% ───────────────────────────────────────────────────────────────
\section{Metode yang Diusulkan}
\label{sec:method}

\subsection{Gambaran Umum Pipeline}

Gambar~\ref{fig:pipeline} mengilustrasikan \textit{pipeline} persiapan
data empat langkah yang diterapkan sebelum pelatihan model.
\textit{Pipeline} ini mengubah dataset mentah yang sangat tidak seimbang
menjadi set pelatihan dengan distribusi kelas yang lebih terkontrol,
memungkinkan aliran gradien yang lebih adil di antara 120 kelas
selama optimasi.

\begin{figure}[htbp]
\centerline{\includegraphics[width=\columnwidth]{fig1.png}}
\caption{Pipeline persiapan data empat langkah. Langkah 1--4
  diterapkan secara berurutan sebelum pelatihan.}
\label{fig:pipeline}
\end{figure}

\subsection{Langkah 1: Stratified Split Per Kelas}

Dataset dibagi menjadi set pelatihan (70\%), validasi (15\%),
dan pengujian (15\%) menggunakan pembagian stratifikasi
\textit{per kelas}.
Hal ini menjamin semua 120 kelas terwakili di setiap subset,
bahkan untuk kelas minoritas yang hanya memiliki 21 gambar,
menghasilkan $\approx$15 gambar pelatihan, 3 gambar validasi, dan
3 gambar pengujian per kelas minoritas sebelum augmentasi.
Pembagian ini menghasilkan 1.789 gambar validasi dan 1.919 gambar
pengujian.

\subsection{Langkah 2: Integrasi Data Variasi}

Sebanyak 1.066 gambar dari \texttt{variations/} ditambahkan
ke set pelatihan 8 kelas konsonan Vokal A yang bersesuaian
(ba, ga, ha, ka, la, nya, ra, ta).
Langkah ini meningkatkan keberagaman gaya tulisan tangan untuk
kelas-kelas tersebut tanpa memerlukan pengumpulan data baru.

\subsection{Langkah 3: Augmentasi Berbasis Preservasi Struktur}

Ini merupakan kontribusi metodologis utama penelitian.
Augmentasi diterapkan \textit{hanya pada kelas minoritas} yang
jumlah gambar pelatihannya berada di bawah ambang batas target
$T = 300$ gambar, yaitu mencakup seluruh 100 kelas non-A.
Gambar-gambar baru dibuat dengan mengambil sampel berulang dari
gambar pelatihan yang ada dan menerapkan transformasi hingga
setiap kelas mencapai $T$ sampel.

Batasan desain yang kritis: \textbf{tidak menggunakan transformasi
geometris}.
Dua observasi dari literatur menjustifikasi ini:
(1) Wisesty et al.~\cite{wisesty2026} secara empiris menunjukkan
bahwa augmentasi menurunkan kinerja pengenalan aksara Jawa,
yang dikaitkan dengan transformasi yang mengganggu struktur
visual karakteristik;
(2) Pasangan ``ha''/``la''~\cite{nindya2025} merupakan cerminan
horizontal yang hampir identik, sehingga penerapan \textit{horizontal
flip} langsung menyebabkan korupsi label lintas kelas.

Tabel~\ref{tab:aug} mencantumkan transformasi yang diizinkan dan
dilarang.
Transformasi yang diizinkan menargetkan \textit{variabilitas
tampilan} (kerapatan tinta, derau pemindaian, ketebalan goresan
pena) tanpa mengubah tata letak spasial karakter.
Setiap gambar yang diaugmentasi dibuat dengan menerapkan 1--3
transformasi yang dipilih secara acak dari daftar yang diizinkan.
Augmentasi dilakukan secara \textit{offline} (disimpan ke disk)
untuk menjamin reprodusibilitas.
Hasil akhir set pelatihan setelah augmentasi berjumlah
\textbf{38.510 gambar} (meningkat $\approx4{,}5\times$ dari 8.534
gambar asli).

\begin{table}[htbp]
\caption{Strategi Augmentasi Berbasis Preservasi Struktur}
\label{tab:aug}
\begin{center}
\renewcommand{\arraystretch}{1.2}
\begin{tabular}{p{1.7cm}p{3.6cm}c}
\toprule
\textbf{Kategori} & \textbf{Transformasi} & \textbf{Status} \\
\midrule
\multirow{5}{1.7cm}{Fotometrik}
  & Jitter kecerahan ($\pm$30\%)                  & \checkmark \\
  & Jitter kontras ($\pm$30\%)                    & \checkmark \\
  & Gaussian blur (radius 0,5--1,5\,px)           & \checkmark \\
  & Gaussian noise ($\sigma$ = 1--5\%)            & \checkmark \\
  & Penajaman ($1{,}2\times$--$2{,}0\times$)      & \checkmark \\
\midrule
\multirow{2}{1.7cm}{Morfologis}
  & Dilasi (MaxFilter $3\!\times\!3$ / $5\!\times\!5$) & \checkmark \\
  & Erosi (MinFilter $3\!\times\!3$ / $5\!\times\!5$)  & \checkmark \\
\midrule
Spasial & Crop acak + resize (maks 10\% tepi) & \checkmark \\
\midrule
\multirow{4}{1.7cm}{Geometris (dilarang)}
  & Rotasi ($>\pm3^{\circ}$)            & $\times$ \\
  & Flip horizontal / vertikal          & $\times$ \\
  & Shear / transformasi afin           & $\times$ \\
  & Elastic distortion                  & $\times$ \\
\bottomrule
\end{tabular}
\end{center}
\end{table}

\subsection{Langkah 4: Fungsi Loss Berbobot}

Setelah augmentasi, bobot kelas dihitung menggunakan rumus:
\begin{equation}
  w_c = \frac{N}{C \cdot n_c}
  \label{eq:weight}
\end{equation}
dengan $N$ adalah total sampel pelatihan, $C = 120$ adalah jumlah
kelas, dan $n_c$ adalah jumlah sampel pelatihan pada kelas $c$.
Formula ini setara dengan \texttt{compute\_class\_weight('balanced')}
pada scikit-learn, yang memberikan penalti kerugian lebih besar
pada kesalahan klasifikasi kelas minoritas.
Bobot ini diterapkan sebagai parameter pada fungsi
\texttt{CrossEntropyLoss} selama pelatihan.

\subsection{Arsitektur Model}

Empat arsitektur CNN dengan bobot pralatih ImageNet-1K digunakan
dalam studi komparatif ini.
Pemilihan arsitektur didasarkan pada relevansi dengan literatur
aksara Jawa yang ada, sebagaimana dirangkum pada Tabel~\ref{tab:arch}.
Kepala klasifikasi setiap arsitektur digantikan dengan lapisan
linear yang memetakan ke 120 kelas.

\begin{table}[htbp]
\caption{Arsitektur CNN yang Dibandingkan}
\label{tab:arch}
\begin{center}
\renewcommand{\arraystretch}{1.2}
\begin{tabular}{lcl}
\toprule
\textbf{Arsitektur} & \textbf{Parameter} & \textbf{Relevansi Literatur} \\
\midrule
EfficientNet-B0~\cite{tan2019efficientnet}    & $\sim$5,3 juta  & -- \\
MobileNetV3-Large~\cite{howard2019mobilenetv3}& $\sim$5,4 juta  & Wisesty et al.~\cite{wisesty2026} \\
ResNet-18~\cite{he2016resnet}                 & $\sim$11,7 juta & Ifriza et al.~\cite{ifriza2026} \\
VGG-16~\cite{simonyan2015vgg}                 & $\sim$138 juta  & Nindya et al.~\cite{nindya2025} \\
\bottomrule
\end{tabular}
\end{center}
\end{table}

\subsection{Strategi Pelatihan}

Semua model mengikuti pendekatan \textit{fine-tuning} dua fase
yang seragam:
\begin{itemize}
  \item \textbf{Fase 1 (10 epoch):}
    Tulang punggung (\textit{backbone}) dibekukan; hanya kepala
    klasifikasi yang dilatih menggunakan optimizer Adam.
    Fase ini membangun inisialisasi kepala yang stabil sebelum
    gradien merambat ke fitur pralatih.
  \item \textbf{Fase 2 (hingga 90 epoch):}
    Seluruh lapisan dibebaskan untuk \textit{fine-tuning} penuh
    dengan learning rate yang lebih rendah.
    ReduceLROnPlateau (faktor 0,5, \textit{patience} 3) menyesuaikan
    laju pembelajaran berdasarkan stagnasi validasi.
    \textit{Early stopping} menghentikan pelatihan ketika validasi
    \textit{loss} tidak membaik selama sejumlah epoch tertentu.
\end{itemize}

Konfigurasi hyperparameter per arsitektur ditunjukkan pada
Tabel~\ref{tab:hyperparam}.
Loss menggunakan \texttt{CrossEntropyLoss} berbobot sesuai
persamaan~\eqref{eq:weight}.
Citra dikonversi ke \textit{grayscale} (3-channel) dan
dinormalisasi dengan statistik ImageNet
($\mu\,{=}\,[0{,}485;\ 0{,}456;\ 0{,}406]$,
 $\sigma\,{=}\,[0{,}229;\ 0{,}224;\ 0{,}225]$).

\begin{table}[htbp]
\caption{Konfigurasi Hyperparameter per Arsitektur}
\label{tab:hyperparam}
\begin{center}
\renewcommand{\arraystretch}{1.2}
\begin{tabular}{lcccc}
\toprule
\textbf{Arsitektur} & \textbf{LR Fase 1} & \textbf{LR Fase 2}
  & \textbf{WD} & \textbf{ES} \\
\midrule
EfficientNet-B0    & $10^{-3}$           & $10^{-4}$          & $10^{-5}$ & 10 \\
MobileNetV3-Large  & $10^{-3}$           & $10^{-4}$          & $10^{-6}$ & 8  \\
ResNet-18          & $2{\times}10^{-3}$  & $2{\times}10^{-4}$ & $10^{-5}$ & 12 \\
VGG-16             & $5{\times}10^{-4}$  & $2{\times}10^{-5}$ & $10^{-4}$ & 15 \\
\bottomrule
\multicolumn{5}{l}{\footnotesize WD = \textit{weight decay}; ES = \textit{patience early stopping}.}
\end{tabular}
\end{center}
\end{table}

% ───────────────────────────────────────────────────────────────
\section{Pengaturan Eksperimen}
\label{sec:setup}

\subsection{Perangkat Keras dan Perangkat Lunak}

Semua eksperimen dilakukan pada sistem dengan GPU NVIDIA GeForce
RTX 3070 Ti (8\,GB VRAM) dan CPU AMD Ryzen 5 3600.
Perangkat lunak: Python 3.10, PyTorch 2.9.1 (CUDA 13.0),
torchvision 0.24.1, scikit-learn.
Seluruh \textit{random seed} ditetapkan pada nilai 42 untuk
reprodusibilitas.
Ukuran \textit{batch}: 32; jumlah \textit{worker} DataLoader: 3.

\subsection{Metrik Evaluasi}

Metrik berikut dilaporkan pada set pengujian (1.919 sampel):
\begin{itemize}
  \item \textbf{Keseluruhan:} Akurasi Top-1, F1-score makro, dan
    F1-score tertimbang (\textit{weighted}).
    F1-score makro dipilih sebagai metrik utama karena
    memberikan bobot yang sama kepada setiap kelas, sehingga
    tidak bias terhadap kelas mayoritas.
  \item \textbf{Per kelompok vokal:} F1-score makro dihitung
    secara terpisah untuk (a) 20 kelas Vokal A (mayoritas,
    $\approx$1.519 sampel uji) dan (b) 100 kelas Vokal non-A
    (minoritas, 400 sampel uji, 4 per kelas).
    Ini merupakan kontribusi analitis utama yang memperlihatkan
    kesenjangan generalisasi model pada kondisi imbalance.
  \item \textbf{Kesenjangan A vs.\ non-A:}
    Selisih absolut F1 antara kelompok mayoritas dan minoritas
    digunakan untuk mengukur keseimbangan kinerja model.
\end{itemize}

\subsection{Catatan Limitasi Pengujian}

Karena set pengujian mengikuti distribusi stratifikasi, setiap
kelas minoritas hanya memiliki \textbf{4 gambar} dalam set uji.
Nilai F1 per kelas pada tingkat granularitas ini memiliki
variansi yang sangat tinggi; satu prediksi yang salah
memengaruhi F1 hingga 25\%.
Oleh karena itu, metrik per kelompok vokal (rata-rata atas 20
kelas) lebih dapat diandalkan dibandingkan nilai per kelas
individual.

% ───────────────────────────────────────────────────────────────
\section{Hasil dan Pembahasan}
\label{sec:results}

\subsection{Perbandingan Kinerja Keseluruhan}

Tabel~\ref{tab:main} merangkum kinerja semua arsitektur pada
set pengujian.
EfficientNet-B0 mencapai akurasi tertinggi 97,08\% dan F1-score
tertimbang 0,9048, sementara ResNet-18 mencapai F1-score makro
tertinggi (0,9180).
MobileNetV3-Large dan VGG-16 menunjukkan kinerja lebih rendah
meskipun masing-masing memiliki 5,4 dan 138 juta parameter.

\begin{table*}[htbp]
\caption{Perbandingan Kinerja Empat Arsitektur pada Set Pengujian 120 Kelas Aksara Jawa}
\label{tab:main}
\begin{center}
\renewcommand{\arraystretch}{1.2}
\begin{tabular}{lccccccc}
\toprule
\textbf{Arsitektur} & \textbf{Parameter}
  & \textbf{Akurasi (\%)} & \textbf{F1 Makro}
  & \textbf{F1 Tertimbang}
  & \textbf{F1 Vokal A} & \textbf{F1 Vokal non-A}
  & \textbf{Epoch} \\
\midrule
EfficientNet-B0    & $\sim$5,3 juta  & \textbf{97,08} & 0,9110 & \textbf{0,9048} & \textbf{0,9902} & 0,8783 & 50 \\
ResNet-18          & $\sim$11,7 juta & 96,93          & \textbf{0,9180} & -- & 0,8972 & \textbf{0,8976} & 36 \\
VGG-16             & $\sim$138 juta  & 93,38          & 0,7867 & --     & 0,9824 & 0,7402 & 36 \\
MobileNetV3-Large  & $\sim$5,4 juta  & 90,83          & 0,6935 & 0,9048 & 0,9714 & 0,6321 & 52 \\
\bottomrule
\multicolumn{8}{l}{\footnotesize \textbf{Cetak tebal} = nilai terbaik per kolom. Epoch = jumlah epoch saat \textit{early stopping} dipicu.}\\
\multicolumn{8}{l}{\footnotesize F1 Tertimbang untuk ResNet-18 dan VGG-16 tidak tersimpan; ``--'' tidak berarti nol.}
\end{tabular}
\end{center}
\end{table*}

\subsection{Analisis Kesenjangan Vokal A vs.\ Vokal non-A}

Tabel~\ref{tab:gap} menyajikan analisis kesenjangan antara kelas
mayoritas dan minoritas, yang merupakan kontribusi analitis utama
penelitian ini.
Hasil ini mengungkap dua pola penting yang tersembunyi oleh
metrik agregat.

\begin{table}[htbp]
\caption{Kesenjangan F1 Makro: Vokal A (Mayoritas) vs.\ Vokal non-A (Minoritas)}
\label{tab:gap}
\begin{center}
\renewcommand{\arraystretch}{1.2}
\begin{tabular}{lccc}
\toprule
\textbf{Arsitektur} & \textbf{F1 Vokal A} & \textbf{F1 non-A}
  & \textbf{Kesenjangan} \\
\midrule
EfficientNet-B0    & 0,9902 & 0,8783 & 0,1119 \\
ResNet-18          & 0,8972 & 0,8976 & \textbf{0,0004} \\
VGG-16             & 0,9824 & 0,7402 & 0,2422 \\
MobileNetV3-Large  & 0,9714 & 0,6321 & 0,3393 \\
\bottomrule
\multicolumn{4}{l}{\footnotesize Kesenjangan = $|$F1 Vokal A $-$ F1 Vokal non-A$|$.
  \textbf{Cetak tebal} = terkecil.}
\end{tabular}
\end{center}
\end{table}

\textbf{Pola pertama: EfficientNet-B0 unggul secara akurasi, tetapi
tidak seimbang.}
Akurasi 97,08\% dan F1 Vokal A 0,9902 menunjukkan bahwa arsitektur
ini sangat baik dalam mengenali kelas mayoritas, tetapi F1 Vokal
non-A 0,8783 dan kesenjangan 0,1119 mengindikasikan bahwa kelas
minoritas masih lebih sering salah diklasifikasikan.

\textbf{Pola kedua: ResNet-18 adalah arsitektur paling seimbang.}
Kesenjangan antara F1 Vokal A (0,8972) dan F1 Vokal non-A (0,8976)
hanya sebesar 0,0004, yang secara praktis dapat diabaikan.
Ini menunjukkan bahwa ResNet-18 dengan residual connection
berhasil mempelajari fitur yang sama baiknya untuk kelas dengan
data melimpah maupun kelas dengan data yang telah diaugmentasi.
Temuan ini tidak terlihat pada studi-studi sebelumnya yang hanya
menggunakan dataset seimbang.

VGG-16 dan MobileNetV3-Large menunjukkan kesenjangan yang lebih
besar (0,2422 dan 0,3393), mengindikasikan bahwa kedua arsitektur
ini lebih rentan terhadap bias mayoritas meskipun \textit{pipeline}
yang sama diterapkan.

\subsection{Analisis Kurva Pelatihan}

Gambar~\ref{fig:curves_eff} dan~\ref{fig:curves_res} menampilkan
kurva pelatihan untuk dua arsitektur terbaik.
Transisi antara Fase 1 (epoch 1--10) dan Fase 2 (epoch 11
dan seterusnya) terlihat jelas sebagai penurunan \textit{loss}
yang tajam dan kenaikan akurasi yang drastis, mengkonfirmasi
efektivitas strategi pelatihan dua fase.

\begin{figure}[htbp]
\centerline{\includegraphics[width=\columnwidth]%
  {../results_efficientnet_b0/training_curves.png}}
\caption{Kurva pelatihan EfficientNet-B0. Transisi Fase 1 ke
  Fase 2 terlihat pada epoch 11.}
\label{fig:curves_eff}
\end{figure}

\begin{figure}[htbp]
\centerline{\includegraphics[width=\columnwidth]%
  {../results_resnet_18/training_curves.png}}
\caption{Kurva pelatihan ResNet-18. Konvergensi yang lebih
  cepat dibandingkan EfficientNet-B0 (36 vs.\ 50 epoch).}
\label{fig:curves_res}
\end{figure}

EfficientNet-B0 berhenti pada epoch ke-50 dengan validation
loss 0,0988 dan validation accuracy 97,60\%.
ResNet-18 berhenti lebih awal (epoch ke-36), mencerminkan
konvergensi yang lebih efisien.
VGG-16 memerlukan waktu komputasi per-epoch yang jauh lebih
lama ($\approx$1.470 detik per epoch pada Fase 2 vs.\
$\approx$120 detik untuk EfficientNet-B0), sehingga tidak
cocok untuk eksperimen iteratif dengan dataset ini.

\subsection{Analisis Matriks Konfusi}

Gambar~\ref{fig:cm_eff} menampilkan matriks konfusi
$120\!\times\!120$ untuk EfficientNet-B0.
Diagonal dominan mengkonfirmasi kinerja yang tinggi secara
keseluruhan, namun blok di luar diagonal yang terkonsentrasi
pada baris-baris Vokal non-A mengindikasikan kesalahan
klasifikasi lintas vokal yang lebih sering terjadi pada
kelas minoritas.

\begin{figure}[htbp]
\centerline{\includegraphics[width=\columnwidth]%
  {../results_efficientnet_b0/confusion_matrix.png}}
\caption{Matriks konfusi EfficientNet-B0 ($120\!\times\!120$).
  Blok bawah-kanan (Vokal non-A) menunjukkan kesalahan
  yang lebih bervariasi dibandingkan blok kiri-atas (Vokal A).}
\label{fig:cm_eff}
\end{figure}

Berdasarkan temuan Nindya et al.~\cite{nindya2025}, pasangan
karakter ``ha''/``la'' (CLR 25,3\%) dan ``sa''/``da'' (CLR 15,8\%)
diprediksi menjadi penyumbang utama kesalahan klasifikasi pada
kelas Vokal A.
Ketiadaan \textit{horizontal flip} dalam strategi augmentasi
yang diusulkan secara langsung bertujuan memitigasi kebingungan
lintas kelas untuk pasangan cerminan seperti ``ha''/``la''.

\subsection{Perbandingan dengan Penelitian Sebelumnya}

Tabel~\ref{tab:comparison} mengontekstualisasikan penelitian ini
dalam literatur pengenalan aksara Jawa yang ada.
Penelitian ini adalah yang pertama mengevaluasi 120 kelas silabik
pada dataset tidak seimbang secara alami, sekaligus yang pertama
melaporkan perbandingan empat arsitektur dengan analisis per
kelompok vokal.

\begin{table}[htbp]
\caption{Perbandingan dengan Penelitian Pengenalan Aksara Jawa}
\label{tab:comparison}
\begin{center}
\renewcommand{\arraystretch}{1.2}
\begin{tabular}{p{2.4cm}cccp{1.6cm}}
\toprule
\textbf{Studi} & \textbf{Kl.} & \textbf{Seimbang?}
  & \textbf{Metrik Terbaik} & \textbf{Arsitektur} \\
\midrule
Susanto et al.~\cite{susantoIJECE2023} & 120 & Ya & Akurasi 97,29\% & Custom CNN \\
Susanto et al.~\cite{susantoIJAI2023}  & 120 & Ya & Akurasi 99,65\% & CNN 12-layer \\
Ifriza et al.~\cite{ifriza2026}        & 20  & Ya & Akurasi 95,53\% & ResNet-18 \\
Wisesty et al.~\cite{wisesty2026}      & 20  & Ya & F1\,=\,0,979   & MobileNetV3 \\
Nindya et al.~\cite{nindya2025}        & 20  & Ya & Akurasi 93,19\% & VGG-16 \\
Faizin et al.~\cite{faizin2025}        & 20  & Tidak & F1 meningkat & SkelBAGAN \\
\midrule
\textbf{Usulan ini} & \textbf{120} & \textbf{Tidak}
  & \textbf{Akurasi 97,08\%; F1 Makro 0,918}
  & \textbf{4 Arsitektur} \\
\bottomrule
\multicolumn{5}{l}{\footnotesize ``Seimbang?'' = apakah dataset pelatihan diseimbangkan secara artifisial.}
\end{tabular}
\end{center}
\end{table}

EfficientNet-B0 yang diusulkan mencapai akurasi sebanding dengan
studi seimbang terdahulu (97,08\% vs.\ 97,29\% Susanto et al.),
meskipun dilatih pada dataset dengan ketidakseimbangan 33,10$\times$.
Hal ini mengkonfirmasi bahwa \textit{pipeline} yang diusulkan
berhasil menjembatani kondisi nyata yang tidak seimbang mendekati
kinerja kondisi seimbang ideal.

% ───────────────────────────────────────────────────────────────
\section{Kesimpulan}
\label{sec:conclusion}

Makalah ini memaparkan \textit{pipeline} data empat langkah dan
studi komparatif empat arsitektur \textit{transfer learning} untuk
pengenalan 120 kelas aksara Jawa tulisan tangan di bawah
ketidakseimbangan kelas yang parah (IR\,=\,33,10$\times$).
Kontribusi metodologis utama adalah strategi augmentasi berbasis
preservasi struktur yang secara eksplisit mengecualikan transformasi
geometris yang terbukti mengubah makna visual aksara~%
\cite{wisesty2026,nindya2025}, dikombinasikan dengan stratified split,
integrasi data variasi, dan \textit{weighted cross-entropy loss}.

Dari evaluasi komparatif diperoleh dua temuan utama.
\textbf{Pertama,} EfficientNet-B0 mencapai akurasi tertinggi
97,08\% dengan F1-score makro 0,9110, sebanding dengan studi
yang menggunakan dataset seimbang~\cite{susantoIJECE2023}.
\textbf{Kedua,} ResNet-18 menunjukkan keseimbangan kinerja
yang luar biasa dengan kesenjangan F1 antara Vokal A dan
Vokal non-A hanya 0,0004 (0,8972 vs.\ 0,8976), menjadikannya
pilihan arsitektur terbaik ketika kesesuaian generalisasi
antar kelas lebih diprioritaskan daripada akurasi agregat.
MobileNetV3-Large dan VGG-16, meskipun relevan dengan literatur
sebelumnya, menunjukkan kesenjangan yang lebih besar (0,3393 dan
0,2422) pada kondisi ketidakseimbangan ini.

Penelitian lanjutan mencakup:
(i) sintesis kelas minoritas berbasis GAN sebagaimana~\cite{faizin2025}
sebagai sumber augmentasi komplementer yang dapat mengurangi
ketergantungan pada augmentasi fotometrik;
(ii) peningkatan jumlah gambar pengujian per kelas minoritas
untuk memperoleh estimasi metrik per kelas yang lebih andal
(saat ini hanya 4 gambar/kelas minoritas);
dan (iii) perluasan ke pengenalan tingkat kata termasuk bentuk
\textit{pasangan} (penggabungan konsonan), yang merupakan batas
terbuka berikutnya sebagaimana diidentifikasi
oleh~\cite{wicaksono2024}.

% ───────────────────────────────────────────────────────────────
\begin{thebibliography}{00}

\bibitem{faizin2025}
M. A. Faizin, N. Suciati, dan C. Fatichah,
``Handling imbalance in Javanese manuscript character dataset
using skeleton-based balancing generative adversarial networks,''
\textit{Jurnal RESTI (Rekayasa Sistem dan Teknologi Informasi)},
vol.~9, no.~4, hal.~820--831, 2025.

\bibitem{he2016resnet}
K. He, X. Zhang, S. Ren, dan J. Sun,
``Deep residual learning for image recognition,''
dalam \textit{Proc. IEEE Conf. Computer Vision and Pattern Recognition
(CVPR)}, 2016, hal.~770--778.

\bibitem{howard2019mobilenetv3}
A. Howard et al.,
``Searching for MobileNetV3,''
dalam \textit{Proc. IEEE/CVF Int. Conf. Computer Vision (ICCV)},
2019, hal.~1314--1324.

\bibitem{ifriza2026}
Y. N. Ifriza et al.,
``Optimizing Javanese script recognition using fine-tuned ResNet-18
and transfer learning,''
\textit{International Journal of Artificial Intelligence (IJ-AI)},
vol.~15, no.~1, hal.~443--453, 2026.

\bibitem{nindya2025}
T. P. Nindya, M. D. Sulistiyo, dan U. N. Wisesty,
``Recognition of Balinese, Javanese, and Sundanese scripts
using CNN with VGG-16 architecture,''
dalam \textit{Proc. 2025 Int. Conf. Data Science and Its Applications
(ICoDSA)}, 2025, hal.~1221--1226.

\bibitem{simonyan2015vgg}
K. Simonyan dan A. Zisserman,
``Very deep convolutional networks for large-scale image
recognition,''
dalam \textit{Proc. Int. Conf. Learning Representations (ICLR)},
2015.

\bibitem{sulistiyo2025}
M. D. Sulistiyo et al.,
``Aksarawa: Comparative study of shallow learning for OCR
of Nusantara scripts,''
\textit{Jurnal RESTI}, vol.~9, no.~6, hal.~1537--1546, 2025.

\bibitem{susantoIJAI2023}
A. Susanto et al.,
``Handwritten Javanese script recognition method based 12-layers
deep convolutional neural network and data augmentation,''
\textit{International Journal of Artificial Intelligence (IJ-AI)},
vol.~12, no.~3, hal.~1448--1458, 2023.

\bibitem{susantoIJECE2023}
A. Susanto, I. U. W. Mulyono, C. A. Sari, E. H. Rachmawanto,
dan D. R. I. M. Setiadi,
``Improved Javanese script recognition using custom model of
convolution neural network,''
\textit{International Journal of Electrical and Computer Engineering
(IJECE)}, vol.~13, no.~6, hal.~6629--6636, 2023.

\bibitem{tan2019efficientnet}
M. Tan dan Q. V. Le,
``EfficientNet: Rethinking model scaling for convolutional neural
networks,''
dalam \textit{Proc. Int. Conf. Machine Learning (ICML)},
2019, hal.~6105--6114.

\bibitem{wicaksono2024}
B. A. Wicaksono, B. S. Hantono, dan T. B. Adji,
``Word segmentation task for Southeast Asian abugida scripts:
A systematic literature review,''
dalam \textit{Proc. 2024 2nd Int. Conf. Technology Innovation and
Its Applications (ICTIIA)}, 2024.

\bibitem{widiarti2026}
R. Widiarti dan F. T. Adji,
``A baseline evaluation of OCR segmentation and classification
methods for printed Javanese script,''
\textit{Engineering, Technology \& Applied Science Research (ETASR)},
vol.~16, no.~1, hal.~31699--31705, 2026.

\bibitem{wisesty2026}
U. N. Wisesty et al.,
``MNetNCR: MobileNet model for efficient traditional Nusantara
script character recognition,''
\textit{International Journal of Artificial Intelligence (IJ-AI)},
vol.~15, no.~2, hal.~1513--1528, 2026.

\end{thebibliography}

\end{document}
